\section{The method III: stochastic reaction--diffusion via the moment hierarchy}
\label{sec:stochastic}

Everything so far is deterministic, but our greatest contribution is primarily in the stochastic domain. The stochastic payload enters through a single structural fact: \emph{for the right noise class, the SPDE's moments obey closed deterministic reaction--diffusion equations} --- the exact family \eqref{eq:family} -- \eqref{eq:coupledrd} of \S\ref{sec:method}--\ref{sec:coupled}.

\subsection{Moment systems for the linear $N$-species model}

\todo{hOOOOOLLy idek}

\paragraph{Setting.}
Let $D\subset\mathbb{R}$ be the periodic interval $\mathbb{R}/L\mathbb{Z}$ or a bounded interval with Neumann conditions, and let $\Delta$ be the corresponding Laplacian on $H=L^2(D)$. Let $(e_m)_{m\ge 0}$ be its real orthonormal eigenbasis, $-\Delta e_m=\lambda_m e_m$, $\lambda_m\ge 0$.
 
The noise has spatial covariance $q:D\times D\to\mathbb{R}$, which is symmetric, continuous and positive semidefinite. By Mercer's theorem, $q=\sum_p \mu_p\, f_p\otimes f_p$ with $(f_p)$ an orthonormal basis of $H$, $\mu_p\ge 0$, and $\sum_p\mu_p<\infty$. Put $\bar q=\sup|q|<\infty$. Let
\[
W_k(t)=\sum_p\sqrt{\mu_p}\,\beta_k^p(t)\,f_p ,
\]
where the $\beta_k^p$ are independent standard Brownian motions, so the channels $k$ are independent with common spatial covariance $q$. Consider
\begin{equation}\label{eq:sys-revised}
du_i=\Big(D_i\Delta u_i+\sum_j R_{ij}u_j\Big)dt+\sum_k\sigma_{ik}\,u_k\,dW_k,
\qquad u_i(0)=u_i^0,
\end{equation}
with $D_i\ge 0$, $R,\sigma$ constant matrices, and $u^0$ independent of $W$ with $\mathbb{E}\|u^0\|_H^4<\infty$.
 
\medskip
\noindent\textbf{Hypothesis (H).} \eqref{eq:sys-revised} has a unique mild solution $u\in C([0,T];H)$ a.s. with $\sup_{t\le T}\mathbb{E}\|u(t)\|_H^4<\infty$.
 
(H) holds under the assumptions above. The map $v\mapsto \sum_k\sigma_{ik}\,v\,dW_k$ is Lipschitz $H\to \mathrm{HS}(Q^{1/2}H,H)$, because $\sum_p\mu_p\|v f_p\|^2=\int q(x,x)v(x)^2dx\le\bar q\|v\|^2$. Existence, uniqueness and moment bounds then follow from the standard theory (Da~Prato--Zabczyk, \emph{Stochastic Equations in Infinite Dimensions}, Thm.~7.2 and its $p$-th moment version).
 
\paragraph{Moment objects.}
Define the Bochner expectations
\[
m_i(t)=\mathbb{E}[u_i(t)]\in L^2(D),\qquad
V_{ij}(t)=\mathbb{E}\big[u_i(t)\otimes u_j(t)\big]\in L^2(D\times D),
\]
so that formally $V_{ij}(t,x,y)=\mathbb{E}[u_i(t,x)\,u_j(t,y)]$. These are well defined because $\|u_i\otimes u_j\|=\|u_i\|\|u_j\|$ is integrable. The pointwise (``same-point'') moment is the restriction $V_{ij}(t,x,x)$, which is \emph{not} an object of this proposition.
 
\begin{proposition}[Closed moment systems]\label{prop:moments-revised}
Under (H):
\begin{enumerate}
\item[\textup{(a)}] \textbf{Mean.} $m=(m_i)$ is the unique mild solution of
\begin{equation}\label{eq:mean-revised}
\partial_t m_i=D_i\Delta m_i+\sum_j R_{ij}m_j .
\end{equation}
\item[\textup{(b)}] \textbf{Two-point function.} $V=(V_{ij})$ is the unique mild solution in $C([0,T];L^2(D\times D)^{N\times N})$ of
\begin{equation}\label{eq:two-point-revised}
\begin{aligned}
\partial_t V_{ij}(x,y)
&=\big(D_i\Delta_x+D_j\Delta_y\big)V_{ij}(x,y)
+\sum_k\big(R_{ik}V_{kj}(x,y)+R_{jk}V_{ik}(x,y)\big)\\
&\quad+\,q(x,y)\sum_k\sigma_{ik}\sigma_{jk}\,V_{kk}(x,y),
\end{aligned}
\end{equation}
with $V_{ij}(0)=\mathbb{E}[u_i^0\otimes u_j^0]$. Here $\Delta_x,\Delta_y$ carry the boundary condition of $\Delta$ in each variable separately; in the Neumann case this is $\partial_{n_x}V_{ij}=\partial_{n_y}V_{ij}=0$, in the weak sense.
\end{enumerate}
In particular, \eqref{eq:two-point-revised} is a closed, linear, deterministic system for $(V_{ij})$: the noise source involves only the two-point functions $V_{kk}(x,y)$ and the covariance $q(x,y)$.
 
\smallskip
\noindent\emph{Status:} \textbf{[P]}, proved below modulo the cited well-posedness (H). The weak-sense Neumann condition on $V$ is also a consequence of the proof (it settles the item previously marked ``used but not proven'').
\end{proposition}

\todo{this is claude, i believe goes without saying}
\begin{proof}
\textbf{Step 1 (mode equations).} Let $u_i^m(t)=\langle u_i(t),e_m\rangle$. Testing the mild formulation against $e_m$, and using that $\Delta$ is self-adjoint with $\Delta e_m=-\lambda_m e_m$, gives
\[
u_i^m(t)=u_i^m(0)+\int_0^t\Big(-D_i\lambda_m u_i^m+\sum_j R_{ij}u_j^m\Big)ds+M_i^m(t),
\]
\[
M_i^m(t)=\sum_k\sigma_{ik}\sum_p\sqrt{\mu_p}\int_0^t\langle e_m,u_k f_p\rangle\,d\beta_k^p .
\]
The $\beta_k^p$ are independent across $(k,p)$, so
\[
d\langle M_i^m,M_j^n\rangle_t=\sum_k\sigma_{ik}\sigma_{jk}\sum_p\mu_p\langle e_m,u_kf_p\rangle\langle e_n,u_kf_p\rangle\,dt .
\]
By Mercer's theorem the series $\sum_p\mu_p f_p(x)f_p(y)$ converges uniformly to $q(x,y)$, and $u_k\in L^2\subset L^1$, so we may interchange sum and integral:
\begin{equation}\label{eq:qv-revised}
d\langle M_i^m,M_j^n\rangle_t=\sum_k\sigma_{ik}\sigma_{jk}\iint e_m(x)e_n(y)\,q(x,y)\,u_k(t,x)u_k(t,y)\,dx\,dy\;dt .
\end{equation}
Also $\sum_p\mu_p\langle e_m,u f_p\rangle^2=\iint e_m(x)e_m(y)\,q(x,y)\,u(x)u(y)\le\bar q\|u\|^2$, so $d\langle M_i^m\rangle\le C\bar q\|u\|^2dt$.
 
\textbf{Step 2 (Itô product and expectation).} By the Itô product rule for continuous semimartingales,
\[
d\big(u_i^mu_j^n\big)=u_i^m\,du_j^n+u_j^n\,du_i^m+d\langle M_i^m,M_j^n\rangle .
\]
The stochastic integrals $\int u_i^m\,dM_j^n$ are true martingales: since $|u_i^m|\le\|u_i\|$ and $d\langle M_j^n\rangle\le C\bar q\|u_j\|^2dt$,
\[
\mathbb{E}\int_0^T(u_i^m)^2\,d\langle M_j^n\rangle\le C\bar q\int_0^T\mathbb{E}\big[\|u_i\|^2\|u_j\|^2\big]dt<\infty
\]
by the fourth-moment bound in (H). Taking expectations, writing $V_{ij}^{mn}=\mathbb{E}[u_i^mu_j^n]=\langle V_{ij},e_m\otimes e_n\rangle$, and applying Fubini in \eqref{eq:qv-revised} (the integrand is dominated by $\bar q\,|e_m(x)e_n(y)|\,\mathbb{E}[|u_k(x)||u_k(y)|]$, which is integrable), we get
\begin{equation}\label{eq:coef-ode-revised}
\frac{d}{dt}V_{ij}^{mn}=-(D_i\lambda_m+D_j\lambda_n)V_{ij}^{mn}+\sum_k\big(R_{ik}V_{kj}^{mn}+R_{jk}V_{ik}^{mn}\big)+\sum_k\sigma_{ik}\sigma_{jk}\,\big\langle q\,V_{kk},\,e_m\otimes e_n\big\rangle .
\end{equation}
The right side depends only on the family $\{V_{kl}\}$, not on any higher moment. This is the closure.
 
\textbf{Step 3 (identification and uniqueness).} Let $X=L^2(D\times D)^{N\times N}$. Let $B$ act on the $(i,j)$ component as $D_i\Delta_x+D_j\Delta_y$ with the stated boundary conditions. It is self-adjoint and nonpositive, diagonal in the basis $e_m\otimes e_n$ with eigenvalue $-(D_i\lambda_m+D_j\lambda_n)$, hence generates a contraction $C_0$-semigroup on $X$; the boundary conditions on $V$ are those of the domain of $B$. Let $\mathcal L$ collect the remaining terms of \eqref{eq:two-point-revised}. It is bounded on $X$:
\[
\|(\mathcal L v)_{ij}\|\le\sum_k\big(|R_{ik}|\|v_{kj}\|+|R_{jk}|\|v_{ik}\|\big)+\bar q\sum_k|\sigma_{ik}\sigma_{jk}|\,\|v_{kk}\| ,
\]
because multiplication by the bounded function $q$ is bounded on $L^2(D\times D)$. So $B+\mathcal L$ generates a $C_0$-semigroup on $X$.
 
The map $t\mapsto V(t)\in X$ is continuous, by dominated convergence from continuity of $u$ in $H$ and the fourth-moment bound. By \eqref{eq:coef-ode-revised}, $V$ is a weak solution of $V'=(B+\mathcal L)V$ tested against the basis $e_m\otimes e_n$, which is dense in $D(B)$ in graph norm. A continuous weak solution coincides with the mild solution (Ball, 1977), and the mild solution is unique. This proves (b).
 
\textbf{Part (a)} is identical with the noise term absent: take expectations in Step~1, using $\mathbb{E}M_i^m=0$, to get $\frac{d}{dt}m_i^m=-D_i\lambda_m m_i^m+\sum_jR_{ij}m_j^m$.
\end{proof}
 
\begin{remark}[Why the pointwise equation does not close]\label{rem:pointwise-fails}
The former Prop.~12 asserted $\partial_t V_{ij}(x,x)=(D_i+D_j)\partial_{xx}V_{ij}(x,x)+\cdots$ for the same-point function. Set $v_{ij}(x)=V_{ij}(x,x)$. The chain rule gives
\[
\partial_{xx}\,v_{ij}=\big(\Delta_x+\Delta_y+2\partial_x\partial_y\big)V_{ij}\Big|_{y=x}.
\]
For $D_i=D_j=D$ the diffusion term in \eqref{eq:two-point-revised}, restricted to the diagonal, is therefore
\[
D\,\partial_{xx}v_{ij}\;-\;2D\,\mathbb{E}\big[\partial_xu_i\,\partial_xu_j\big](x),
\]
which contains the gradient moment $\mathbb{E}[\partial_xu_i\,\partial_xu_j]$. That moment is not a function of $v_{ij}$; its own evolution involves higher derivatives of $V$, so the pointwise hierarchy does not close. For $D_i\ne D_j$ the diagonal diffusion term cannot even be written as a multiple of $\partial_{xx}v_{ij}$ plus a gradient term.
 
\emph{Deterministic counterexample ($\sigma=0$, $N=1$, $R=0$).} Take $u_t=Du_{xx}$ and $u_0=2+\cos x$, so $V(0,x,y)=u_0(x)u_0(y)$. At $x=\pi/2$ we have $u_0=2$, $u_{0,x}=-1$, $u_{0,xx}=0$. Then
\[
\partial_t V(x,x)\big|_{t=0}=2D\,u_0u_{0,xx}=0,
\qquad
D\,\partial_{xx}\big[V(x,x)\big]\big|_{t=0}=2D\big(u_{0,x}^2+u_0u_{0,xx}\big)=2D .
\]
The same-point equation fails even for heat flow.
 
\emph{Stationary check.} If $V_{ii}(x,y)=c(x-y)$ is homogeneous, then $v_{ii}$ is constant, so the pointwise equation predicts no diffusive damping of the variance. The two-point equation instead gives a diffusive contribution $2D\,c''(0)=-2D\,\mathbb{E}[(\partial_xu_i)^2]<0$ on the diagonal, as it should.
\end{remark}
 
\begin{remark}[Cases where the diagonal does close]
If $D_i=0$ for all $i$ (no spatial coupling), then \eqref{eq:two-point-revised} contains no spatial derivatives, so the diagonal restriction (with $q(x,x)$ in the source) closes and reduces to the per-cell finite-dimensional It\^o computation. The failure is entirely due to the diffusion term, which couples a point to its neighbours. The mean system \eqref{eq:mean-revised} closes for any noise, because the equation is linear in $u$.
\end{remark}
 
\begin{remark}[Consequences for OTS--NIDC \textup{(status [O])}]
\eqref{eq:two-point-revised} is a $2$-dimensional linear reaction--diffusion system in $(x,y)$ with a potential-like coupling $q(x,y)$. This has three consequences.
\begin{itemize}
\item The scalar analysis of \S1 applies to \eqref{eq:two-point-revised} only after extension to two space dimensions. The squared Laplacian $(D_i\Delta_x+D_j\Delta_y)^2$ contains a mixed term $2D_iD_j\,\partial_{xxyy}$, which the one-dimensional cancellation at $\Delta t_{\mathrm{opt}}$ does not remove; an additional cross-derivative correction is required.
\item $q(x,y)$ is a variable coefficient, so $D_i\Delta$ and the reaction term do not commute; the binomial identity behind Eq.~(2) fails and Group~2 and Group~3 pick up derivatives of $q$.
\item The cost scales as $N_x^2$ for the two-point function, so the efficiency comparison against Euler--Maruyama must be restated.
\end{itemize}
For white noise ($q=\delta(x-y)$) the source is $\delta(x-y)\sum_k\sigma_{ik}\sigma_{jk}V_{kk}(x,x)$; the equation is plausible in the distributional sense in one dimension, but this proof does not cover it (it needs $\bar q<\infty$), and $V$ then has a gradient jump on the diagonal, which affects the Taylor-based order claims. This is open.
\end{remark}
 
\begin{remark}[Notation]
The former text wrote $M_i$ in \eqref{eq:mean-revised}; use $m_i$ for means and reserve $M_k$ for nothing. The old Prop.~11 is subsumed by (a) for $k=1$; it is retracted for $k\ge 2$.
\end{remark}

\todo{decide which proof is actually worth to keep}
\begin{proof}[Alternative proof of Proposition~11]
Write $\mathcal D$ for the spatial domain and $H=L^2(\mathcal D)$.
All expectations and integrals below are justified by hypothesis~(H),
the boundedness of $q$, and the Cauchy--Schwarz inequality.
We derive the equations in weak form, so no pointwise differentiability
of the random field is assumed.

\paragraph{Step 1: The mean equation.}
Fix a species $i$ and a smooth test function $\varphi$ satisfying the
periodic or homogeneous Neumann boundary condition. Testing (18)
against $\varphi$ gives
\begin{align*}
\langle u_i(t),\varphi\rangle
={}&\langle u_i(0),\varphi\rangle
 +\int_0^t\left[D_i\langle u_i(s),\Delta\varphi\rangle
       +\sum_kR_{ik}\langle u_k(s),\varphi\rangle\right]ds\\
&+\sum_{k,p}\sigma_{ik}\sqrt{\mu_p}
   \int_0^t\langle u_k(s)f_p,\varphi\rangle\,d\beta_{kp}(s).
\end{align*}
The last line is an integrable, mean-zero stochastic integral:
indeed, its second moment is bounded by a constant multiple of
$\int_0^t\mathbb E\|u(s)\|_H^2\,ds$, since
$\sum_p\mu_p\|u_k f_p\|_H^2
=\int_{\mathcal D}q(x,x)|u_k(x)|^2\,dx
\leq\bar q\|u_k\|_H^2$.
Taking expectations and setting $m_i=\mathbb E[u_i]$ yields
\[
\langle m_i(t),\varphi\rangle
=\langle m_i(0),\varphi\rangle
 +\int_0^t\left[D_i\langle m_i(s),\Delta\varphi\rangle
       +\sum_kR_{ik}\langle m_k(s),\varphi\rangle\right]ds.
\]
This is the weak form of (19).

\paragraph{Step 2: The two-point product.}
Fix species $i,j$ and smooth test functions $\varphi,\psi$ as above.
Introduce the real-valued processes
\[
X(t)=\langle u_i(t),\varphi\rangle,
\qquad Y(t)=\langle u_j(t),\psi\rangle.
\]
They satisfy ordinary real-valued stochastic integral equations.
For any increments the algebraic identity
\[
\Delta(XY)=X\,\Delta Y+Y\,\Delta X+\Delta X\,\Delta Y
\]
is exact. The finite-dimensional It\^o product rule is the rigorous
continuous-time version of this identity: the product of the two
noise increments contributes their covariance, while drift--noise
and drift--drift products vanish at first order.

For each channel $k$ and covariance mode $p$, the noise coefficients
of $X$ and $Y$ are, respectively,
\[
a_{kp}=\sigma_{ik}\sqrt{\mu_p}\langle u_k f_p,\varphi\rangle,
\qquad
b_{kp}=\sigma_{jk}\sqrt{\mu_p}\langle u_k f_p,\psi\rangle.
\]
Independence of the $\beta_{kp}$ implies that the covariance term
in the expected product increment is
$\mathbb E\sum_{k,p}a_{kp}b_{kp}\,dt$.
The expansion $q(x,y)=\sum_p\mu_p f_p(x)f_p(y)$ therefore gives
\begin{align*}
\sum_{k,p}a_{kp}b_{kp}
&=\sum_k\sigma_{ik}\sigma_{jk}
  \sum_p\mu_p
   \left(\int_{\mathcal D}u_k(x)f_p(x)\varphi(x)\,dx\right)
   \left(\int_{\mathcal D}u_k(y)f_p(y)\psi(y)\,dy\right)\\
&=\sum_k\sigma_{ik}\sigma_{jk}
  \iint_{\mathcal D^2}q(x,y)u_k(x)u_k(y)
                         \varphi(x)\psi(y)\,dx\,dy.
\end{align*}
The interchange of summation and integration is justified by the
covariance-operator bound and $\bar q<\infty$; equivalently it follows
by taking limits of finite covariance-mode sums in $L^1$.

\paragraph{Step 3: Taking expectations.}
The deterministic part of the product rule contributes
\begin{align*}
& D_i\,\mathbb E\!ig[\langle u_i,\Delta\varphi\rangle
                         \langle u_j,\psi\rangle\big]
 +D_j\,\mathbb E\!\big[\langle u_i,\varphi\rangle
                         \langle u_j,\Delta\psi\rangle\big]\\
&\quad+\sum_k R_{ik}\,\mathbb E\!\big[
         \langle u_k,\varphi\rangle\langle u_j,\psi\rangle\big]
 +\sum_k R_{jk}\,\mathbb E\!\big[
         \langle u_i,\varphi\rangle\langle u_k,\psi\rangle\big].
\end{align*}
The remaining stochastic integrals have zero expectation.
For example, their square-integrability follows from (H), because
$|X|^2\leq\|\varphi\|_H^2\|u_i\|_H^2$ and the variance rate of the
noise in $Y$ is bounded by
$C\bar q\|\psi\|_H^2\sum_k\|u_k\|_H^2$.
Thus their expected quadratic integrals are controlled by
$C\int_0^T\mathbb E\|u(s)\|_{H^N}^4\,ds<\infty$.

By definition of $V_{ij}=\mathbb E[u_i\otimes u_j]$, combining these
terms gives the integral identity
\begin{align*}
\langle V_{ij}(t),\varphi\otimes\psi\rangle
={}&\langle V_{ij}(0),\varphi\otimes\psi\rangle\\
&+\int_0^t\Big\{
  D_i\langle V_{ij},\Delta\varphi\otimes\psi\rangle
 +D_j\langle V_{ij},\varphi\otimes\Delta\psi\rangle\\
&\hspace{2.3em}+\sum_k\big(
  R_{ik}\langle V_{kj},\varphi\otimes\psi\rangle
 +R_{jk}\langle V_{ik},\varphi\otimes\psi\rangle\big)\\
&\hspace{2.3em}+\sum_k\sigma_{ik}\sigma_{jk}
   \langle qV_{kk},\varphi\otimes\psi\rangle\Big\}\,ds.
\end{align*}
Here each integrand is evaluated at $s$.
Finite sums of product test functions $\varphi\otimes\psi$ are dense
in $L^2(\mathcal D^2)$ and form a core for the product Laplacian
with the stated boundary conditions. Hence this identity is the weak
form of (20), with the separate spatial operators
$D_i\Delta_x+D_j\Delta_y$.

\paragraph{Step 4: Closure and uniqueness.}
The right-hand side of (20) uses only the collection
$\{V_{k\ell}\}_{k,\ell=1}^N$: the noise term involves $V_{kk}$ rather
than a third- or fourth-order moment. Thus the two-point system closes.
To see uniqueness, let $Z$ be the difference of two mild solutions
with identical initial data. The diffusion operator
$B Z=(D_i\Delta_x Z_{ij}+D_j\Delta_y Z_{ij})_{ij}$
generates a contraction semigroup $S(t)$ on
$\mathcal X=L^2(\mathcal D^2)^{N\times N}$.
The remaining reaction-and-noise operator $L$ is bounded on
$\mathcal X$, since $R,\sigma$ are finite matrices and
$\|qZ_{kk}\|_{L^2}\leq\bar q\|Z_{kk}\|_{L^2}$.
The mild integral equation therefore implies
\[
\|Z(t)\|_{\mathcal X}
\leq\|L\|\int_0^t\|Z(s)\|_{\mathcal X}\,ds.
\]
Gronwall's inequality gives $Z=0$. Existence of the mild solution
follows from the bounded-perturbation theorem for this same
contraction semigroup, and the weak identity above identifies the
moment field with that unique solution. The identical deterministic
argument applies to the mean system (19). This proves (a) and (b).
\end{proof}

Proposition~\ref{prop:closure} makes \S\ref{sec:method} apply verbatim to every moment equation. This formulation buys us a couple nice formulations:

\begin{enumerate}
\item \textbf{One step for the whole hierarchy} (Cor.~\ref{cor:indep}): $\dto = \dx^2/6D$ governs every $M_k$ simultaneously; $\sigma$ and $k$ appear only in the correction coefficients \eqref{eq:C1} --- e.g.\ the mean ($\alpha= \frac12\sigma^2$) and second moment ($\alpha=\sigma^2$) are advanced by the same step, different corrections. Since It\^o drift is necessarily a defect and not a structural change (the noise-induced reaction term is explicitly computable) it slots into $C^{(1)}$ exactly as any reaction term would. 

\item \textbf{Weak, not strong, by construction:} The method never samples $\dot W$, so the Davie--Gaines strong-order-$1/2$ ceiling \citep{davie2001} governs a problem we do not solve. The delivered object is fourth-order convergence of the computed moment fields to the true moments of the closed system \tagV (scalar, Table~\ref{tab:scalar}); conversion to statements about the SPDE law rides entirely on Proposition~\ref{prop:closure}'s exact closure.
  
\item \textbf{Stability transfers:} the growing-moment case is exactly the $\eta\ge0$ regime of Prop.~\ref{prop:stability} --- the correction marginally \emph{shrinks} the CFL window and is comfortably stable at $\dto$ \tagV. \todo{unsure of this}
\end{enumerate}

\emph{The competitor, stated fairly}

Euler--Maruyama on \eqref{eq:spde} estimates $M_k$ from $M$ paths with error
$O(\dt) + O(M^{-1/2})$ per point; our implementation (\texttt{baselines.py},
\texttt{euler\_maruyama\_moment}) exists for exactly this comparison
\tagS\footnote{The EM baseline is implemented and its noise scaling
($\sigma u\sqrt{\dt/\dx}\,Z$ per cell per step) is the standard space-time white
noise discretization; the head-to-head convergence figure against it is not yet
in the results ledger --- it is scheduled with the benchmarking gate of
\S\ref{sec:open}, per the project's own decision to hold benchmarking until
scoping stabilized.}. The structural comparison, however, is already decisive
and does not need the figure: for the \emph{moments alone}, the pathwise route
pays $M$ path-solves plus Monte Carlo variance to approximate what the moment
route computes deterministically --- and the moment route's solver is
fourth-order where EM is weak-first. The pathwise route is irreplaceable when
paths are the product (extremes, hitting, filtering); for moment delivery under
exact closure it is dominated.

\subsection{Where the noise model ends}

Three extensions are visible from inside the current theory, all open:
\begin{itemize}
\item \textbf{Colored noise} (trace-class $Q$): the closure proof goes through
with $\sigma^2 \to$ the covariance trace structure; the M/V systems acquire
nonlocal reaction kernels. The OTS cancellation, being diffusive-local, should
survive with kernel-aware corrections \tagO.
\item \textbf{Noise discretization interplay:} how Wiener-increment sampling on
the grid (piecewise-constant vs.\ linear-in-time noise,
\citep{petersson2022}) interacts with the NIDC cancellation --- for the moment
route there is no sampled noise, which is precisely the advantage; the question
becomes how to \emph{validate} moment solvers against pathwise references
fairly \tagO.
\item \textbf{Boundary conditions for $V$:} Neumann data on $u$ induce
$\partial_n V_{ij}=0$ on $\partial\Omega\times\Omega \cup \Omega\times
\partial\Omega$; used but not proven in the coupled harness \tagO.
\end{itemize}
